%=============================================================================
% APPENDIX V: EXTENDED PHENOMENOLOGY AND NUMERICAL METHODS
% Wide binaries, external field effect, FEM implementation
%=============================================================================

\section{Extended Phenomenology and Numerical Methods}\label{app:extended-phenom}

This appendix addresses three areas that complete the DFD phenomenological 
framework: the external field effect (EFE), wide binary predictions, and 
numerical implementation via finite element methods.

%-----------------------------------------------------------------------------
\subsection{The External Field Effect (EFE)}\label{app:efe}
%-----------------------------------------------------------------------------

In nonlinear theories like DFD, the internal dynamics of a subsystem depend 
on its external gravitational environment. This \emph{external field effect} 
(EFE) arises from the nonlinearity of the field equation.

\subsubsection{Physical Origin}

The DFD field equation
\begin{equation}
\nabla \cdot \bigl[\mu(|\nabla\psi|/a_\star)\nabla\psi\bigr] = -\frac{8\pi G}{c^2}\rho
\end{equation}
is nonlinear in $\nabla\psi$. For a subsystem (e.g., a dwarf galaxy) embedded 
in an external field (e.g., a host galaxy), the total gradient is:
\begin{equation}
|\nabla\psi_{\mathrm{tot}}| = |\nabla\psi_{\mathrm{int}} + \nabla\psi_{\mathrm{ext}}|.
\end{equation}

When $|\nabla\psi_{\mathrm{ext}}| \gg a_\star/c^2$ but $|\nabla\psi_{\mathrm{int}}| \ll a_\star/c^2$, 
the total gradient may exceed the crossover scale even if internal accelerations 
are in the deep-field regime. This ``Newtonianizes'' the internal dynamics.

\subsubsection{Quantitative Formulation}

Define the dimensionless acceleration ratios:
\begin{align}
x_{\mathrm{int}} &= \frac{c^2|\nabla\psi_{\mathrm{int}}|}{2a_0}, \\
x_{\mathrm{ext}} &= \frac{c^2|\nabla\psi_{\mathrm{ext}}|}{2a_0}, \\
x_{\mathrm{tot}} &= \frac{c^2|\nabla\psi_{\mathrm{int}} + \nabla\psi_{\mathrm{ext}}|}{2a_0}.
\end{align}

The effective $\mu$-function argument becomes $x_{\mathrm{tot}}$, not $x_{\mathrm{int}}$:
\begin{equation}
\mu_{\mathrm{eff}} = \mu(x_{\mathrm{tot}}) = \frac{x_{\mathrm{tot}}}{1 + x_{\mathrm{tot}}}.
\end{equation}

\paragraph{Limiting cases.}
\begin{itemize}
\item \textbf{Isolated system} ($x_{\mathrm{ext}} \to 0$): $\mu_{\mathrm{eff}} = \mu(x_{\mathrm{int}})$, 
standard DFD dynamics.
\item \textbf{Strong external field} ($x_{\mathrm{ext}} \gg 1$, $x_{\mathrm{ext}} \gg x_{\mathrm{int}}$): 
$\mu_{\mathrm{eff}} \approx 1$, Newtonian dynamics restored.
\item \textbf{Aligned fields}: Maximum enhancement when $\nabla\psi_{\mathrm{int}} \parallel \nabla\psi_{\mathrm{ext}}$.
\item \textbf{Opposed fields}: Partial cancellation possible.
\end{itemize}

\subsubsection{Observational Signatures}

\begin{table}[h]
\centering
\caption{External field effect predictions for Milky Way satellites.}\label{tab:efe-predictions}
\begin{tabular}{lccc}
\toprule
Satellite & $g_{\mathrm{ext}}$ (m/s$^2$) & $x_{\mathrm{ext}}$ & EFE suppression \\
\midrule
Fornax & $2 \times 10^{-11}$ & 0.17 & Mild (15\%) \\
Sculptor & $3 \times 10^{-11}$ & 0.25 & Moderate (20\%) \\
Draco & $5 \times 10^{-11}$ & 0.42 & Significant (30\%) \\
Crater II & $1 \times 10^{-11}$ & 0.08 & Weak (8\%) \\
\bottomrule
\end{tabular}
\end{table}

The EFE predicts that satellites at smaller galactocentric radii (higher 
$g_{\mathrm{ext}}$) show \emph{less} enhanced dynamics than isolated dwarfs 
with similar internal properties.

\paragraph{Falsification criterion.}
If dwarf satellites uniformly show enhanced dynamics independent of their 
position relative to the Milky Way, the EFE mechanism (and hence DFD's 
nonlinear structure) would be falsified.

%-----------------------------------------------------------------------------
\subsection{Wide Binary Predictions}\label{app:wide-binaries}
%-----------------------------------------------------------------------------

Wide stellar binaries with separations $s \gtrsim 5000$ AU probe the 
low-acceleration regime where DFD deviates from Newtonian gravity.

\subsubsection{The Crossover Scale}

For a binary with total mass $M$ and separation $s$, the internal acceleration is:
\begin{equation}
a_{\mathrm{int}} = \frac{GM}{s^2}.
\end{equation}

The crossover to deep-field behavior occurs when $a_{\mathrm{int}} \sim a_0$:
\begin{equation}
s_{\mathrm{cross}} = \sqrt{\frac{GM}{a_0}} \approx 7000\,\mathrm{AU} \times \left(\frac{M}{M_\odot}\right)^{1/2}.
\end{equation}

For solar-mass binaries, $s_{\mathrm{cross}} \approx 7000$ AU.

\subsubsection{Predicted Velocity Anomaly}

In the deep-field regime ($s \gg s_{\mathrm{cross}}$), the orbital velocity is enhanced:
\begin{equation}
v_{\mathrm{DFD}} = (GMa_0)^{1/4} = v_{\mathrm{Newton}} \times \left(\frac{s}{s_{\mathrm{cross}}}\right)^{1/2}.
\end{equation}

The velocity ratio relative to Newtonian prediction:
\begin{equation}
\frac{v_{\mathrm{DFD}}}{v_{\mathrm{Newton}}} = \sqrt{\frac{a_{\mathrm{Newton}}}{a_0} + 1} \approx 
\begin{cases}
1 & s \ll s_{\mathrm{cross}}, \\
\sqrt{s/s_{\mathrm{cross}}} & s \gg s_{\mathrm{cross}}.
\end{cases}
\end{equation}

\begin{table}[h]
\centering
\caption{DFD predictions for wide binary velocity anomalies.}\label{tab:wide-binary}
\begin{tabular}{lccc}
\toprule
Separation (AU) & $a_{\mathrm{int}}/a_0$ & $v_{\mathrm{DFD}}/v_{\mathrm{Newton}}$ & Observable effect \\
\midrule
1000 & 50 & 1.01 & Negligible \\
3000 & 5.6 & 1.08 & 8\% enhancement \\
7000 & 1.0 & 1.22 & 22\% enhancement \\
10000 & 0.5 & 1.37 & 37\% enhancement \\
20000 & 0.13 & 1.73 & 73\% enhancement \\
\bottomrule
\end{tabular}
\end{table}

\subsubsection{GAIA DR3 Constraints}

Recent analyses of GAIA DR3 wide binary data show conflicting results:
\begin{itemize}
\item Some analyses report enhanced relative velocities consistent with 
MOND-like dynamics at $s > 5000$ AU~\cite{Chae2023}.
\item Other analyses find no significant deviation from Newtonian 
predictions~\cite{Banik2024}.
\end{itemize}

\paragraph{DFD interpretation.}
The EFE complicates wide binary tests: binaries in regions of higher 
galactic acceleration ($g_{\mathrm{gal}} \gtrsim a_0$) are partially 
Newtonianized. A definitive test requires:
\begin{itemize}
\item Selection of binaries in low-$g_{\mathrm{gal}}$ environments.
\item Proper treatment of projection effects and orbital phase.
\item Statistical comparison with DFD predictions including EFE.
\end{itemize}

\paragraph{Falsification criterion.}
If wide binaries in isolated, low-acceleration environments show 
strictly Newtonian dynamics at $s > 10000$ AU, DFD's deep-field 
prediction would be falsified.

%-----------------------------------------------------------------------------
\subsection{Finite Element Implementation}\label{app:fem}
%-----------------------------------------------------------------------------

The DFD field equation is directly implementable via finite element methods 
(FEM). We outline the key elements for numerical solution.

\subsubsection{Weak Form for FEM}

The weak formulation~\eqref{eq:weak-form-U} translates directly to FEM assembly:
\begin{equation}
\sum_e \int_{\Omega_e} \mu(|\nabla\psi_h|)\nabla\psi_h \cdot \nabla v_h\,dx 
= \sum_e \int_{\Omega_e} f\,v_h\,dx,
\end{equation}
where $\psi_h, v_h$ are finite element approximations on mesh elements $\Omega_e$.

\subsubsection{Newton Iteration for Nonlinearity}

The nonlinear system is solved via Newton iteration. The Jacobian matrix is:
\begin{equation}
J_{ij}(\nabla\psi) = \mu(|\nabla\psi|)\delta_{ij} + \mu'(|\nabla\psi|)\frac{\partial_i\psi\,\partial_j\psi}{|\nabla\psi|}.
\end{equation}

For $\mu(s) = s/(1+s)$:
\begin{equation}
\mu'(s) = \frac{1}{(1+s)^2}.
\end{equation}

\paragraph{Regularization at small gradients.}
At $|\nabla\psi| \to 0$, the Jacobian may become ill-conditioned. A standard 
remedy is regularization:
\begin{equation}
|\nabla\psi| \to \sqrt{|\nabla\psi|^2 + \epsilon^2},
\end{equation}
with $\epsilon \sim 10^{-10}$ in dimensionless units.

\subsubsection{Mesh Refinement Strategy}

The deep-field regime features steep gradients near sources. Adaptive mesh 
refinement (AMR) is recommended:
\begin{itemize}
\item Refine where $|\nabla\psi|$ changes rapidly (gradient indicator).
\item Refine near crossover radius $r_\star = \sqrt{GM/a_\star}$.
\item Use logarithmic radial spacing for exterior domains.
\end{itemize}

\subsubsection{Boundary Conditions}

\paragraph{Dirichlet (fixed $\psi$).}
\begin{equation}
\psi|_{\Gamma_D} = \psi_D.
\end{equation}
Used for outer boundaries with known asymptotic value.

\paragraph{Neumann (fixed flux).}
\begin{equation}
\mu(|\nabla\psi|)\nabla\psi \cdot \hat{n}|_{\Gamma_N} = g_N.
\end{equation}
Used for symmetry planes or specified matter flux.

\paragraph{Robin (mixed).}
\begin{equation}
\mu(|\nabla\psi|)\nabla\psi \cdot \hat{n} + \kappa(\psi - \psi_\infty) = 0.
\end{equation}
Used for approximate radiation conditions at finite boundaries.

\subsubsection{Convergence Verification}

For code verification, use the analytic deep-field solution:
\begin{equation}
\psi(r) = \psi_0 - B\ln(r/r_0), \quad B = \frac{2}{c^2}\sqrt{GMa_\star}.
\end{equation}

Richardson extrapolation on mesh sequences should yield:
\begin{equation}
\|\psi_h - \psi_{\mathrm{exact}}\|_{L^2} = \mathcal{O}(h^{p+1}),
\end{equation}
where $p$ is the polynomial order of the elements.

\begin{tcolorbox}[colback=blue!5!white, colframe=blue!50!black, title=FEM Implementation Checklist]
\begin{enumerate}[nosep]
\item Assemble weak form with $\mu(|\nabla\psi|)\nabla\psi$ flux
\item Newton iteration with analytic Jacobian
\item Regularize $|\nabla\psi|$ at small values
\item Adaptive mesh refinement near crossover
\item Verify against analytic deep-field solution
\item Richardson extrapolation for convergence rate
\end{enumerate}
\end{tcolorbox}

%-----------------------------------------------------------------------------
\subsection{Matter Power Spectrum from $\psi$-Screen}\label{app:power-spectrum}
%-----------------------------------------------------------------------------

The $\psi$-screen formalism (Section~\ref{sec:psi-tomography}) predicts 
modifications to the matter power spectrum $P(k)$.

\subsubsection{Scale-Dependent $\psi$ Perturbations}

Density perturbations $\delta\rho$ source $\delta\psi$ via the linearized 
field equation:
\begin{equation}
\nabla^2\delta\psi = -\frac{8\pi G}{c^2}\delta\rho.
\end{equation}

In Fourier space:
\begin{equation}
\delta\tilde{\psi}(k) = \frac{8\pi G}{c^2 k^2}\delta\tilde{\rho}(k).
\end{equation}

The $\psi$-perturbation power spectrum is:
\begin{equation}
P_\psi(k) = \left(\frac{8\pi G}{c^2 k^2}\right)^2 P_\rho(k).
\end{equation}

\subsubsection{Observational Signatures}

The $\psi$-screen affects:
\begin{itemize}
\item \textbf{CMB lensing:} Modified convergence $\kappa$ from $\psi$-gradients.
\item \textbf{Galaxy clustering:} Scale-dependent bias from $\psi$-density correlation.
\item \textbf{Weak lensing:} Modified shear-density relation.
\end{itemize}

These effects are degenerate with dark matter at leading order but 
distinguishable through their scale dependence and cross-correlations.

%-----------------------------------------------------------------------------
\subsection{Summary}\label{app:extended-summary}
%-----------------------------------------------------------------------------

\begin{tcolorbox}[colback=green!5!white, colframe=green!60!black, title=Extended Phenomenology Summary]
\textbf{External Field Effect:}
\begin{itemize}[nosep]
\item Nonlinear $\mu$-function causes environmental dependence
\item Satellites in strong external fields are Newtonianized
\item Testable via dwarf galaxy velocity dispersions vs. position
\end{itemize}

\textbf{Wide Binaries:}
\begin{itemize}[nosep]
\item Crossover at $s_{\mathrm{cross}} \sim 7000$ AU for solar-mass binaries
\item 20--70\% velocity enhancement predicted for $s > 10000$ AU
\item EFE complicates interpretation; requires low-$g_{\mathrm{gal}}$ samples
\end{itemize}

\textbf{Numerical Methods:}
\begin{itemize}[nosep]
\item Standard FEM with Newton iteration for nonlinearity
\item Regularization needed at small $|\nabla\psi|$
\item Adaptive mesh refinement near crossover scale
\item Verification against analytic deep-field solution
\end{itemize}
\end{tcolorbox}
